\section{Evolución y Resultados}

% --- Ponente 3 ---

\begin{frame}{Evolución Arquitectónica}
\begin{itemize}
    \item \textbf{Primera generación:} CNNs --- capturan patrones \textit{locales} eficientemente.
    \item \textbf{Segunda generación:} \textbf{Vision Transformers (ViT)} --- entienden el \textit{contexto global} mediante mecanismos de \textbf{atención}.
    \item La atención permite ponderar regiones distantes de la imagen para una codificación más eficiente.
\end{itemize}

% INSERTAR FIGURA 2 AQUI
% \begin{figure}
%     \centering
%     \includegraphics[width=0.8\textwidth]{images/evolucion_arquitecturas}
%     \caption{De CNNs a Vision Transformers en compresión neuronal.}
% \end{figure}
\end{frame}

\begin{frame}{Resultados vs. JPEG}
\begin{alertblock}{BD-Rate: 65\% de ahorro}
Los modelos NIC actuales ahorran un \textbf{65\% de bits} respecto a JPEG manteniendo la misma calidad visual en el dataset Kodak.
\end{alertblock}
\begin{itemize}
    \item Métrica \textbf{BD-Rate}: mide el ahorro porcentual de bits a calidad equivalente.
    \item Superan también a HEVC/Intra en múltiples benchmarks.
    \item Rendimiento consistente en distintos tipos de contenido.
\end{itemize}

% INSERTAR FIGURA 4 AQUI (Curva R-D)
% \begin{figure}
%     \centering
%     \includegraphics[width=0.8\textwidth]{images/curva_rd}
%     \caption{Curva Tasa-Distorsión: NIC vs. códecs clásicos.}
% \end{figure}
\end{frame}

\begin{frame}{Ventajas Perceptivas y Adaptabilidad}
\begin{itemize}
    \item \textbf{Fine-tuning} para dominios específicos:
    \begin{itemize}
        \item Radiografías médicas.
        \item Imágenes satelitales.
        \item Contenido artístico.
    \end{itemize}
    \item Uso de métricas perceptuales como \textbf{MS-SSIM} para optimizar para el ojo humano.
    \item Degradación \textbf{suave y gradual}, sin artefactos de bloque.
\end{itemize}

% INSERTAR FIGURA 5 AQUI
% \begin{figure}
%     \centering
%     \includegraphics[width=0.8\textwidth]{images/comparacion_perceptual}
%     \caption{Comparación perceptual: NIC vs. JPEG a misma tasa de bits.}
% \end{figure}
\end{frame}
